\documentclass[a4paper,11pt]{ctexart}
\usepackage{xcolor}
\usepackage{array}
\usepackage{booktabs}
\usepackage{hyperref}
\hypersetup{colorlinks=true,linkcolor=blue!70!black}
\linespread{1.35}

\newcommand{\draft}{\textcolor{orange!80!black}{\textbf{【待写】}}}

\title{\textcolor{blue!70!black}{\Huge\textbf{土木工程模块 · 课程大纲}}}
\author{}
\date{}

\begin{document}
\maketitle
\thispagestyle{empty}

\section*{第零层 · 入门故事}

\begin{enumerate}
\item[\textbf{第1课}] \textbf{桥为什么不会塌？——从赵州桥到港珠澳大桥} \draft \\
\textbf{内容：} 桁架结构、拱形受力、悬索桥原理。北京在地探索：卢沟桥和玉蜓桥的结构对比。
\item[\textbf{第2课}] \textbf{北京最高的楼怎么建起来的？——地基与抗震} \draft \\
\textbf{内容：} 国贸三期330米的秘密——深桩基础、隔震支座。北京在地震带上——摩天大楼如何"抗震"而不是"抗倒"？
\item[\textbf{第3课}] \textbf{大兴机场——全球最大的单体航站楼是怎么建成的？} \draft \\
\textbf{内容：} 新机场的巨型钢网壳屋顶。用什么样的起重机和安装方案把1.4万吨钢结构吊装到位的？
\end{enumerate}

\section*{深入课程}

\begin{enumerate}
\item[\textbf{第1课}] \textbf{结构力学基础} \draft \\
\textbf{内容：} 荷载分类（恒载/活载/风载/地震荷载）→计算简图→轴力/剪力/弯矩图→结构位移计算。情景：为什么悬索桥的主缆可以做到跨过2000米的江面？
\item[\textbf{第2课}] \textbf{混凝土与钢筋混凝土} \draft \\
\textbf{内容：} 混凝土的抗压和钢筋的抗拉如何互补。配筋率、锚固长度、保护层厚度。情景：为什么混凝土开裂是"正常"的？钢筋生锈为什么是"致命"的？
\item[\textbf{第3课}] \textbf{钢结构——从铆钉到高强螺栓} \draft \\
\textbf{内容：} 钢材等级（Q235, Q345, Q390）、连接方式（焊接/螺栓/铆钉）、稳定问题（压杆稳定、局部屈曲）。情景：为什么2018年北京某工地的钢结构坍塌了？
\item[\textbf{第4课}] \textbf{岩土工程与地基处理} \draft \\
\textbf{内容：} 土压力理论、地基承载力、桩基础设计（摩擦桩/端承桩）。情景：北京CBD的"中国尊"（中信大厦528米）——它的地基打了多深？
\item[\textbf{第5课}] \textbf{建筑抗震设计} \draft \\
\textbf{内容：} 抗震设防烈度、结构周期与谱加速度的关系、隔震与消能减震技术。情景：2008年汶川地震的建筑倒塌模式——什么结构"柱强梁弱"还是"梁强柱弱"？
\item[\textbf{第6课}] \textbf{施工技术与管理} \draft \\
\textbf{内容：} BIM（建筑信息模型）、施工组织设计、网络计划（关键路径CPM）。情景：北京大兴机场从开工到投运只用了4年——是如何做到的？
\end{enumerate}

\end{document}
